% !TeX encoding = utf-8
% !TeX program = pdflatex
%% Tsinghua Science and Technology (Tsinghua Sci Technol) - LaTeX Template

\documentclass{TST}

% Исправляем отступы для биографий
\makeatletter
% % Увеличиваем ширину области под фотографию
\def\@TSTBIOhangwidth{1.3in}  % было 1.1in
% Увеличиваем глубину (высоту) области под фотографию  
\def\@TSTBIOhangdepth{1.5in}  % было 1.25in
% Увеличиваем размер фотографии
\def\@TSTBIOphotowidth{1.0in}  % было 1in
\def\@TSTBIOphotodepth{1.25in}  % было 1.25in
\makeatother


\newcommand{\blu}[1]{\textcolor{blue}{#1}}

\TSTsetup{
type      = {Research Article},
doi       = {xxxx},
title     = {Heterogeneous Computational Platform for Parallel Graph Processing},
author    = {Vladimir M. Kozlov},
runauthor = {V. M. Kozlov},
abstract  = {
Efficient graph partitioning is a key factor affecting the performance and scalability of distributed graph databases. The placement of graph vertices across multiple storage nodes determines the amount of inter-shard communication required during graph traversal and shortest-path queries. Modern partitoning algorithms can be roughly devided in two groups: offline and online, in which former take significantly more time than the latter but provide better results. This paper presents a hybrid graph partitioning approach that combines the online partitioning of Fennel algorithm and refintment approach used in METIS framework. The proposed method first constructs an initial balanced partitioning using Fennel and subsequently applies KL-based local optimization to reduce the number of edges crossing shard boundaries while preserving partition balance. This combination leverages the scalability of streaming partitioning and the optimization capabilities of iterative local refinement.
},
keywords  = {computing system, graph partitioning},
copyright = {The Author(s)},
received  = 2026-06-14,
revised   = 2026-06-xx,
accepted  = 2026-06-xx,
}

\address{Alexey Yu. Popov, Stanislav V. Ibragimov, Maxim P. Kalitventsev, Mikhail A. Geine, Jiajian Li are with Bauman Moscow State Technical University, Moscow, 105005, Russian Federation. E-mail: alexpopov@bmstu.ru}
\address{Andrey A. Belevantsev is with the Federal State Budgetary Institution of Science Institute for System Programming of the Russian Academy of Sciences (ISP RAS), Moscow, 109004, Russia. Phone: +7(495) 912-07-54 (ext. 400), E-mail: info-isp@ispras.ru}

% Добавьте в преамбулу документа настройки для листингов
\usepackage{listings}
\usepackage{xcolor}
\usepackage{caption}
\usepackage{pgfplots}
\pgfplotsset{compat=1.18} % Укажите версию совместимости
\usepackage{tikz}
\usepackage{graphicx}

\lstset{
   basicstyle=\fontsize{4}{5}\selectfont\ttfamily, % Экстремально маленький
   numbers=none,
   frame=none,
   backgroundcolor=\color{white},
   breaklines=true,
   xleftmargin=0pt,
   xrightmargin=0pt,
   aboveskip=1pt,
   belowskip=1pt,
   % numbers=left,
   % numberstyle=\tiny\color{black},   
   columns=fullflexible,
}

% Настройка подписей листингов - выравнивание по левому краю
\captionsetup[lstlisting]{
    format=hang, % Выравнивание по левому краю
    justification=raggedright, % Выравнивание по левому краю
    singlelinecheck=false, % Не центрировать короткие подписи
    labelfont=bf, % Жирный шрифт для метки
    labelsep=space, % Разделитель между метки и текстом
    skip=5pt, % Отступ от листинга
    belowskip=5pt, % Отступ снизу
    font=small % Размер шрифта подписи
}

\begin{document}

\maketitle

% % Сноска о финансировании на первой странице
% \footnotetext{Funding: This research was supported by the Russian Science Foundation grant No. 25-11-90002, \url{https://rscf.ru/project/25-11-90002/}.}

\section{Introduction}
Graph-structured data have become a fundamental component of modern information systems, including social networks, recommendation engines, knowledge graphs, transportation networks, and communication infrastructures. As the size of these graphs continues to grow, graph databases increasingly rely on distributed architectures in which graph data are partitioned across multiple storage nodes. While distribution enables horizontal scalability, it also introduces the challenge of efficiently assigning graph vertices to partitions.

The quality of graph partitioning has a direct impact on the performance of distributed graph databases. During graph traversals and shortest-path queries, communication between storage nodes is required whenever adjacent vertices reside in different partitions. Such inter-partition communication incurs significantly higher latency than local memory access and can become the dominant factor limiting system performance. Consequently, an effective partitioning strategy should simultaneously minimize the graph cut while maintaining a balanced distribution of data across storage nodes.

Numerous graph partitioning techniques have been proposed to address this problem. Some graph databases, such as ArangoDB \cite{arangodb_github}, prioritize balanced data distribution using hash-based partitioning rather than optimizing partition quality. Traditional partitioning frameworks, such as METIS \cite{MetisOG}, produce high-quality partitions by exploiting the structural properties of the graph. However, these methods are primarily designed for static graphs and often require global repartitioning, making them computationally expensive for large-scale or dynamically evolving datasets. In contrast, streaming partitioning algorithms, such as Fennel \cite{Fennel}, provide scalable and efficient vertex placement as the graph is constructed, balancing partition sizes while reducing the expected graph cut. Despite their computational efficiency, streaming methods are inherently greedy and may produce suboptimal partitions that cannot be improved without additional refinement. Furthermore, because they make decisions without knowledge of the complete graph structure, their partition quality may deteriorate when vertices are processed in an unfavorable order.

This work investigates the potential of a hybrid graph partitioning approach that combines the scalability of the Fennel streaming algorithm with the local optimization capabilities of the Kernighan--Lin (KL) algorithm. In the proposed method, Fennel is employed to construct an initial balanced partitioning of the graph, after which KL refinement is applied to reduce the number of inter-partition edges while preserving partition balance. The proposed combination aims to achieve a favorable trade-off between partitioning efficiency and partition quality, making it suitable for large distributed graph databases where both scalability and communication cost are critical considerations.

The hypothesis of this work is that Fennel provides an initial partitioning of sufficient quality such that a small number of KL refinement iterations can achieve higher partition quality while substantially reducing computational cost compared to the complete METIS partitioning process.


\section{Architectural Design Methodology}
One of the significant problems in creating a computing complex that performs graph storage and processing is applying a systems approach to its design, i.e., determining the most rational option for dividing software and hardware across system hierarchy levels. This means that early development stages must define the number and functional purpose of system elements for each hierarchical level. The following eight levels were identified (Figure \ref{fig:system_hierarhy}):

\begin{figure*}[h!t]
    \centering
    \begin{minipage}{0.3\textwidth}
        \centering
        \includegraphics[width=\linewidth]{system_hierarhy_en.png}
        \caption{System hierarchy levels}\label{fig:system_hierarhy}
    \end{minipage}
    % \hfill
    \hspace{70pt}
    \begin{minipage}{0.33\textwidth}
        \centering
        \includegraphics[width=\linewidth]{functional_hierarhy_en.png}
        \caption{Software functional abstraction levels}\label{fig:functional_hierarhy}
    \end{minipage}
    % \captionof{figure}{Hierarchical structure of the Teragraph computing complex}\label{fig:teragraph_hierarhy}
\end{figure*}

\begin{itemize}
    \item \textit{Key-Value Memory}. In our system, vertex and edge sets are organized as key-value pairs stored using contemporary semiconductor memory devices based on DDR4 DRAM technology. Since multiple processing elements access a shared storage medium, mechanisms for partitioning the unified address space among them and coordinating request ordering must be established. Memory access operates in batch mode, requiring the formation of key-value packets for efficient data transfer. Beyond the keys and values that represent graph vertices and edges, the memory also contains index structures that enable efficient data access and support graph processing algorithms.
    
    \item \textit{Parallel Memory Access}. A significant problem in multithreaded processing is memory bandwidth. To increase it, memory is divided into independent storage devices, with access through hardware queues and arbitration devices.
    
    \item \textit{Associative-Address Translation}. Information extraction from sets inherently assumes associative access, whereas physical storage media operate on address-based principles. The fundamental challenge of mapping arbitrary keys to specific storage addresses represents a significant technical problem that admits multiple solution approaches. Currently, B+ trees serve as the most practical index structure for this translation layer, offering a balanced compromise between search efficiency, storage overhead, and update performance. However, the search for optimal alternatives continues, with promising candidates including prefix trees (tries), hash-based structures, and hybrid designs that may offer better performance characteristics for specific graph workloads or key distributions.

    \item \textit{Discrete Mathematics Instructions}. This level requires defining hardware composition and efficient implementation of primitive set operations underlying the discrete mathematics instruction set DISC. This involves creating specialized microprocessor devices and microcode for control units.
    
    \item \textit{Heterogeneous Core}. Set processing or their combinations relies not only on discrete mathematics commands but also on arithmetic-logical processing. Universal microprocessors are also necessary for creating familiar programming models. Therefore, universal riscv64im microprocessor (CPE) and lnh64 microprocessor with DISC instruction set (SPE) are combined into a unified heterogeneous core.
    
    \item \textit{Network Interaction}. Since the computational resources of individual heterogeneous cores are inherently limited, effective system scaling requires high-speed inter-core communication networks. Such networks must provide both high throughput and low latency, particularly for the small data bursts characteristic of graph processing workloads. These communication capabilities enable two fundamental scaling strategies: (1) horizontal scaling through coordinated execution across large numbers of parallel heterogeneous cores, and (2) pipelined processing where sequential data streams are processed across multiple specialized cores in assembly-line fashion.
    
    \item \textit{Compute Node}. Multiple heterogeneous cores must be combined into a compute node performing: initialization, task launch and execution management, hardware and software resource access, program debugging, remote user network interaction, etc.
    
    \item \textit{Computing Complex}. Multiple nodes can be connected, achieving greater horizontal scaling, increasing heterogeneous core count, and expanding available memory for set storage. The compute complex level applies IaaS (Infrastructure as a Service) cloud approach for resource management, fully supported by developed system libraries and hardware resources.
\end{itemize}

It is worth noting the positive technical effect achieved through the proposed structuring in Figure \ref{fig:system_hierarhy}: requirements for each level's function set are formed, enabling further independent modernization of each level's structure and implementation. This also enables planning research and development work.

For heterogeneous cores and compute nodes, programming support is provided through C/C++ libraries and language extensions. Four levels of software functional abstraction were defined (Figure \ref{fig:functional_hierarhy}):

\begin{itemize}
    \item \textit{Discrete mathematics instruction level} — a foundational set of 21 DISC machine instructions for basic set operations at the heterogeneous core level;
    
    \item \textit{Set description, extraction, and modification level} — language constructs that extend set processing capabilities to support more complex iterative operations. This includes the definition of composite keys and values as structured data types (using C structures, Python classes, or even JSON-like descriptions), as well as mechanisms for sequential iteration over set elements;
    
    \item \textit{Graph operation level} — programmatically callable functions for graph description and processing algorithms;
    
    \item \textit{Task management and data distribution level} — software tools for controlling and managing distributed processing resources to ensure scalability and enable pipelined execution.
\end{itemize}

This hierarchy of software-accessible functions into layered libraries addresses several important objectives: it ensures relative independence for future development and modernization; enables consistent representation of application tasks across different abstraction levels; and allows the formation of abstract models corresponding to each level of the system hierarchy. An important outcome of the structural analysis of the computing complex's hardware and software components should be the development of generalized system and functional models, ultimately forming an open standard for associative computations analogous to OpenCL. Consequently, this article aims not only to present the architecture and operating principles of a computing system for set and graph processing based on discrete mathematics instructions, but also to propose a conceptual framework for standardizing the abstraction levels of such heterogeneous associative systems.

\section{Conclusion}
This work investigated a hybrid graph partitioning approach that combines the scalability of the Fennel streaming partitioning algorithm with the local optimization capabilities of the Kernighan--Lin refinement algorithm. The proposed method replaces the computationally expensive coarsening and uncoarsening stages of the METIS framework with online partitioning using Fennel, while periodically applying KL refinement to improve partition quality.

The experimental evaluation demonstrated that the proposed approach consistently reduced the graph cut compared to the partitioning produced by Fennel alone. Furthermore, the refinement process typically converged after only two iterations for each pair of partitions and never required more than three iterations in the evaluated experiments. These observations support the hypothesis that the partitioning generated by Fennel is of sufficiently high quality for efficient KL refinement, allowing the computational cost of refinement to remain low while improving partition quality.

Although the current evaluation was performed on synthetic graphs of moderate size, the results indicate that the proposed hybrid approach represents a practical compromise between partition quality and computational efficiency. It is therefore well suited for distributed graph databases and other systems that require scalable online graph partitioning while minimizing inter-partition communication.


\section*{Acknowledgments}
This research was supported by the Russian Science Foundation grant No. 25-11-90002, \url{https://rscf.ru/project/25-11-90002/}

\section*{Declaration of Competing Interest}
The authors declare that they have no known competing financial interests or personal relationships that could have appeared to influence the work reported in this paper.

\bibliographystyle{TSTbib}
\bibliography{export-data}

\section*{Author Biography}

\begin{biography}[popov]{Alexey Yu. Popov} 
Professor, Doctor of Science at Bauman Moscow State Technical University, Russia. Research interests: heterogeneous computing systems, graph processing architectures, discrete mathematics acceleration, and high-performance computing. Published over 60 papers; leads Teragraph project.
\end{biography}

\vspace{0.5cm}

\begin{biography}[ibragimov]{Stanislav V. Ibragimov}
Senior Lecturer at the Department of Computer Systems and Networks, Bauman Moscow State Technical University. His research focuses on network architectures for high-performance computing, distributed systems, and hardware-software co-design.
\end{biography}

\vspace{0.5cm}

\begin{biography}[kalitventsev]{Maksim P. Kalitventsev}
M.Sc., PhD student at Bauman Moscow State Technical University. Specializes in Network-on-Chip (NoC) architectures, simulation modeling, and hardware verification (UVM). Research focuses on communication fabrics for large-scale graph processing systems.
\end{biography}

\vspace{0.5cm}

\begin{biography}[geine]{Mikhail A. Geine}
M.Sc., PhD student and Research Engineer at Bauman Moscow State Technical University. His research interests lie in graph query languages, specialized DBMS, and their application for efficient data management and analytics on heterogeneous supercomputer architectures.
\end{biography}

\vspace{0.5cm}

\begin{biography}[li]{Jiajian Li}
Junior Research Fellow at Bauman Moscow State Technical University, specializing in graph-based knowledge structures and artificial intelligence. His research focuses on the development of semantic and event-driven models for intelligent systems. 
\end{biography}

\vspace{0.5cm}

\begin{biography}[belevancev]{Andrey A. Belevantsev}
Andrey Belevantsev is a Corresponding Member of the Russian Academy of Sciences (RAS, 2025), Doctor of Physical and Mathematical Sciences, and Professor of RAS. His research focuses on program optimizations, static analysis, safe compilation, and parallel programming. Author of 82 scientific works and 11 patents, and holds professor positions at Moscow State University and HSE University.
\end{biography}


\end{document}